% Pursuing Stacks (cotes 134-2 to 134-8): archive cote 134-5, site page 24 (PDF page 25), typescript page 211.
% typescript. Draft compared: ps4.tex:1090-1146.
% AI-transcribed by Claude Opus 5.5 (claude-opus-5-5, adaptive thinking, effort high), Message Batches API, batch ps2 (msgbatch_01NVwr1Yuo8qJf5rykeGBXPK), 2026-11-12; not yet checked by eye.
\scanpage{134-5/24}
\mspage{211}
the final (one point) topos, \over{\unsure{where}}{whence} D(T) is just the homotopy category
(Hot), there \xout{is} \inserted[typed]{must be} of course a vast variety of ways of defining D(T) in term\unsure{s}
of model categories, and I would like to review some which seem significan\unsure{t.}
% NOTE: the right edge of the scan cuts off the ends of several lines
In fact, one cannot help but looking at two mutually dual groups of model
categories, giving rise to \inserted[typed]{(at least)} two definitely non equivalent derived categori\unsure{es}
D(T) and D'(T) say, which,in case T=I\^{}, would correspond to D(I) and
D(I\over{\illegible{1}}{${}^{o}$}). A typical model category for the former \xout{wouldxbe} is made up with
Illusie's semisimplicial sheaves; a typical model category for \xout{thexlatter}
D'(T), on the other hand, should be made up with 1-stacks on T (in Giraud\hand{'}
s sense), for a suitable notion of weak equivalence between these.

Maybe the most natural models of all, in this context, for ``relativ\unsure{e}
homotopy types over the topos T'', should be \uline{\unsure{topoI|topoi}} \inserted{X} over T (generalizing \inserted{the}
categories over an ``indexing category'' I). The only trouble with this
point of view though is that the best we can hope for, in terms of Illu-%
sie's D(T) say, is that a topos X over T gives rise to a proobject of
D(T), which needs not come from an object of D(T) itself i.e. is not neces-%
sarily \xout{locally} ``essentially constant.'' This brings us back to the simpler
\inserted[typed]{and still more basic} question of associating a pro-homotopy type, namely a pro-object of (Hot),
to any topos X - namely \inserted[typed]{back} to the \hand{\v{ }}Cech-Verdier-Artin-Mazur construction.
This has been handled so far using the semisimplicial models for (Hot),
I suspect though that using (Cat) as a modelizer will give a more elegant
treatment, as already suggested earlier. Whichever way we choose to get
the basic functor
\[ \text{Topoi} \longrightarrow \text{Pro(Hot)} \quad , \]
this functor will allow us, given a basic localizer $\underline{W}$ , to define
$\underline{W}$-equivalences between topoi\hand{\unsure{,}} as maps \xout{suchxthat} which become isomorphisms
under the composition of the basic functor above, and the canonical
functor
\[ \text{Pro(Hot}\over{\illegible{1}}{(}\underline{W})) \]
\[ \text{Pro(Hot)} \longrightarrow \text{Pro(Hot}\over{\illegible{1}}{(}\underline{W})) \]
deduced from the localization functor \xout{Hot}
\[ \text{(Hot)} \longrightarrow \text{(Hot(}\underline{W})) \ . \]
(It turns out, using Quillen's theory, that usual weak equivalence is
indeed the finest of all possible basic localizers $\underline{W}$ , hence Hot($\underline{W}$) is
indeed a lo\over{\illegible{1}}{\unsure{ca}}lization of (Hot).)

Maybe it is not too unreasonable to expect that all, or most homotopy
constructions, involving a topos T, can be expressed \xout{by} replacing T by
it's image in Pro(Hot) , or in Pro(Hot($\underline{W}$)) if the constructions are rela-%
tive to a given basic localizer $\underline{W}$ . The very first example one would like

\margin{24}
